\chapter{Binding and Precedence}
\label{ch:binding-precedence}
\label{app:precedence}

The official grammar is primary.  The following precedence convention permits
parentheses to be omitted in expository formulas after unique readability has
been established.

\begin{center}
\begin{tabularx}{\textwidth}{@{}c>{\raggedright\arraybackslash}p{0.38\textwidth}>{\raggedright\arraybackslash}X@{}}
\toprule
Rank & Symbols & Association \\
\midrule
1 & function application, predicate application & determined by arity \\
2 & \(=\) & nonassociative \\
3 & \(\neg,\Box,\Diamond,\forall x,\exists x\) & prefix; governs the immediately following formula \\
4 & \(\land\) & left associative \\
5 & \(\lor\) & left associative \\
6 & \(\to\) & right associative \\
7 & \(\leftrightarrow\) & nonassociative \\
\bottomrule
\end{tabularx}
\end{center}

Thus \(\neg P\land Q\to R\) abbreviates
\(((\neg P)\land Q)\to R\), while
\(P\to Q\to R\) abbreviates \(P\to(Q\to R)\).  Parentheses must be restored
whenever omission could make the intended parse uncertain.  In particular,
\(\forall x\,P(x)\to Q(x)\) abbreviates
\((\forall x\,P(x))\to Q(x)\); write
\(\forall x\,(P(x)\to Q(x))\) when the implication belongs to the
quantifier's scope.

